\section*{Bloque B. Aplicaciones cuantitativas y contrastes econométricos}

\noindent\textit{Preguntas 31 a 65. Evalúan la capacidad de cálculo analítico de estimadores, varianzas residuales, pruebas de significancia, diagnósticos de heterocedasticidad (White, Breusch-Pagan), autocorrelación (Durbin-Watson, Prais-Winsten), sesgo de variable omitida, variables instrumentales (2SLS) y pruebas de especificación.}

\q{En una regresión lineal simple $Y_i = \beta_0 + \beta_1 X_i + u_i$ con $n = 50$ observaciones, se tiene: $\sum (X_i - \bar{X})^2 = 400$, $\sum (Y_i - \bar{Y})^2 = 900$ y $\sum (X_i - \bar{X})(Y_i - \bar{Y}) = 600$. ¿Cuál es el valor estimado de la pendiente MCO $\hat{\beta}_1$?}
\begin{opts}
    \item $\hat{\beta}_1 = 1.50$
    \item $\hat{\beta}_1 = 1.00$
    \item $\hat{\beta}_1 = 0.67$
    \item $\hat{\beta}_1 = 2.25$
\end{opts}

\q{Para los datos de la pregunta anterior, determine el coeficiente de correlación lineal de Pearson $r$ y el coeficiente de determinación $R^2$ del modelo:}
\begin{opts}
    \item $r = 0.75$ y $R^2 = 0.56$
    \item $r = 0.50$ y $R^2 = 0.25$
    \item $r = 0.80$ y $R^2 = 0.64$
    \item $r = 1.00$ y $R^2 = 1.00$
\end{opts}

\q{En un modelo de regresión con $n = 45$ observaciones y $k = 5$ parámetros estimados, la Suma de Cuadrados Residuales es $SCR = 800$. ¿Cuál es la varianza residual estimada ($s^2$) y el error estándar de la regresión ($s$)?}
\begin{opts}
    \item $s^2 = 17.78$ y $s = 4.22$
    \item $s^2 = 16$ y $s = 4.00$
    \item $s^2 = 20$ y $s = 4.47$
    \item $s^2 = 25$ y $s = 5.00$
\end{opts}

\q{En una regresión de demanda de agua, el coeficiente estimado de la variable precio es $\hat{\beta}_1 = -0.45$ con un error estándar de $EE(\hat{\beta}_1) = 0.15$. El tamaño muestral es $n = 120$ y el modelo tiene $k = 4$ parámetros. Para contrastar $H_0: \beta_1 = 0$ frente a $H_1: \beta_1 \neq 0$ con un nivel de significancia del $5\%$ (valor crítico bilateral $t_{0.025, 116} \approx 1.98$):}
\begin{opts}
    \item El estadístico es $t = -3.00$, por lo que se rechaza $H_0$ al $5\%$.
    \item El estadístico es $t = -1.50$, por lo que no se rechaza $H_0$ al $5\%$.
    \item El estadístico es $t = -0.067$, resultado no significativo al $5\%$.
    \item El estadístico es $t = -6.75$, resultado no significativo al $5\%$.
\end{opts}

\q{Un modelo de regresión múltiple con $k = 4$ parámetros (3 variables explicativas más constante) y $n = 64$ observaciones arroja un coeficiente de determinación de $R^2 = 0.60$. ¿Cuál es el valor del estadístico $F$ de significancia global de la regresión?}
\begin{opts}
    \item $F = 30.0$
    \item $F = 15.0$
    \item $F = 45.0$
    \item $F = 20.0$
\end{opts}

\q{Se contrastan 2 restricciones lineales simultáneas ($q = 2$) en un modelo no restringido con $k = 6$ parámetros y $n = 106$ observaciones. La suma de cuadrados residuales no restringida es $SCR_{NR} = 500$ y la restringida es $SCR_R = 560$. El valor del estadístico $F$ es:}
\begin{opts}
    \item $F = 8.0$
    \item $F = 15.0$
    \item $F = 6.0$
    \item $F = 12.0$
\end{opts}

\q{La estimación de una función de demanda residencial de agua arroja el modelo: $\ln(Q) = 2.4 - 0.35 \ln(P) + 0.40 \ln(Y)$, donde $Q$ es el consumo mensual en $m^3$, $P$ es la tarifa por $m^3$ y $Y$ es el ingreso familiar. Si la tarifa de agua sube en un $10\%$, manteniendo el ingreso constante:}
\begin{opts}
    \item el consumo de agua se reduce en $3.5$ metros cúbicos exactos.
    \item el consumo de agua se reduce en aproximadamente un $3.5\%$.
    \item el consumo de agua aumenta en $4.0\%$ por efecto sustitución.
    \item la elasticidad precio de la demanda cambia a $-0.035$.
\end{opts}

\q{En una ecuación semilogarítmica de costos operativos de EPS: $\ln(OPEX) = 8.2 + 0.04 \, Conexiones$, donde $Conexiones$ está expresada en miles de conexiones activas. Si una EPS añade 5 mil conexiones adicionales a su red:}
\begin{opts}
    \item sus costos operativos aumentan en aproximadamente $20\%$.
    \item sus costos operativos aumentan en un $4\%$ acumulado.
    \item la elasticidad costo de las conexiones es igual a $0.04$.
    \item sus costos operativos aumentan en 200 soles mensuales.
\end{opts}

\q{En un modelo lineal-log $Y = 50 + 200 \ln(X)$, si la variable $X$ se incrementa en un $1\%$:}
\begin{opts}
    \item la variable dependiente $Y$ aumenta en $2$ unidades.
    \item la variable dependiente $Y$ aumenta en un $1\%$ exacto.
    \item la variable dependiente $Y$ aumenta en $200$ unidades.
    \item la variable dependiente $Y$ se duplica proporcionalmente.
\end{opts}

\q{Al analizar la multicolinealidad entre dos variables explicativas en un modelo de costos de saneamiento, la regresión auxiliar de una de ellas sobre las demás arroja un coeficiente de determinación de $R_j^2 = 0.80$. ¿Cuál es el Factor de Inflación de la Varianza ($VIF$) correspondiente?}
\begin{opts}
    \item $VIF = 1.25$
    \item $VIF = 5.0$
    \item $VIF = 10.0$
    \item $VIF = 4.0$
\end{opts}

\q{En un modelo con dos regresores y constante, si la matriz $(X'X)^{-1}$ tiene en su diagonal principal los elementos $[0.5, 0.04, 0.09]$ y la varianza residual estimada es $s^2 = 25$, ¿cuál es el error estándar del segundo estimador de pendiente $EE(\hat{\beta}_2)$?}
\begin{opts}
    \item $EE(\hat{\beta}_2) = 1.50$
    \item $EE(\hat{\beta}_2) = 5.00$
    \item $EE(\hat{\beta}_2) = 2.25$
    \item $EE(\hat{\beta}_2) = 0.75$
\end{opts}

\q{Para contrastar heterocedasticidad mediante la prueba de White en un modelo con 2 regresores sin interacciones (la regresión auxiliar incluye constante, $X_1, X_2, X_1^2, X_2^2, X_1 X_2$, total 5 regresores excluyendo constante), se tiene $n = 100$ y el $R^2$ auxiliar es $R_{aux}^2 = 0.15$. El estadístico $LM$ de White y su distribución son:}
\begin{opts}
    \item $LM = 1.50 \sim t_{95}$
    \item $LM = 15.0 \sim \chi^2_5$
    \item $LM = 30.0 \sim F_{5, 95}$
    \item $LM = 7.5 \sim \chi^2_2$
\end{opts}

\q{En una prueba de Breusch-Pagan con $m = 3$ variables explicativas en la ecuación de la varianza y $n = 80$ observaciones, el $R^2$ de la regresión auxiliar de residuos al cuadrado es $R_{BP}^2 = 0.10$. El estadístico $LM = n R_{BP}^2$ es:}
\begin{opts}
    \item $LM = 2.4$
    \item $LM = 0.8$
    \item $LM = 26.67$
    \item $LM = 8.0$
\end{opts}

\q{Si se conoce que la varianza del error en un modelo de corte transversal es proporcional al cuadrado de una variable $X_i$ ($Var(u_i|X) = \sigma^2 X_i^2$), la transformación adecuada de Mínimos Cuadrados Ponderados (WLS) sobre la ecuación $Y_i = \beta_0 + \beta_1 X_i + u_i$ exige estimar:}
\begin{opts}
    \item $\frac{Y_i}{X_i^2} = \frac{\beta_0}{X_i^2} + \frac{\beta_1}{X_i} + \frac{u_i}{X_i^2}$
    \item $\ln(Y_i) = \beta_0 + \beta_1 \ln(X_i) + u_i$
    \item $\frac{Y_i}{X_i} = \beta_0 \left(\frac{1}{X_i}\right) + \beta_1 + \frac{u_i}{X_i}$
    \item $Y_i X_i = \beta_0 X_i + \beta_1 X_i^2 + u_i X_i$
\end{opts}

\q{En una regresión de series de tiempo trimestrales de consumo de agua con $n = 100$ observaciones, el coeficiente de autocorrelación muestral de primer orden de los residuos es $\hat{\rho} = 0.60$. ¿Cuál es el valor aproximado del estadístico de Durbin-Watson $d$?}
\begin{opts}
    \item $d \approx 2.00$
    \item $d \approx 1.20$
    \item $d \approx 3.20$
    \item $d \approx 0.80$
\end{opts}

\q{Si los límites críticos de Durbin-Watson para $n = 100$ y $k' = 2$ regresores al nivel del $5\%$ son $d_L = 1.63$ y $d_U = 1.72$, y el valor calculado es $d = 0.80$, la conclusión estadística es:}
\begin{opts}
    \item no existe evidencia de autocorrelación serial en el error.
    \item la prueba resulta no concluyente por zona de indecisión.
    \item existe evidencia de autocorrelación serial negativa.
    \item existe evidencia de autocorrelación serial positiva.
\end{opts}

\q{Si en un modelo de series de tiempo el estadístico de Durbin-Watson resulta ser $d = 1.30$, ¿cuál es el valor estimado del parámetro autorregresivo $\hat{\rho}$ para aplicar la transformación de Cochrane-Orcutt?}
\begin{opts}
    \item $\hat{\rho} \approx 0.35$
    \item $\hat{\rho} \approx -0.35$
    \item $\hat{\rho} \approx 0.15$
    \item $\hat{\rho} \approx 0.70$
\end{opts}

\q{Bajo la transformación cuasi-diferenciada de Prais-Winsten con parámetro $\hat{\rho} = 0.35$, la variable dependiente transformada para el primer periodo ($t=1$) se define como:}
\begin{opts}
    \item $Y_1^* = (1 - 0.35^2) \, Y_1$
    \item $Y_1^* = \sqrt{1 - (0.35)^2} \, Y_1$
    \item $Y_1^* = (1 - 0.35^2)^{-1/2} \, Y_1$
    \item $Y_1^* = Y_1 - 0.35 Y_0$
\end{opts}

\q{El verdadero modelo poblacional es $Y = 10 + 4 X_1 + 2 X_2 + u$, pero el investigador omite por error la variable $X_2$ y estima por MCO la regresión simple $Y = \tilde{\beta}_0 + \tilde{\beta}_1 X_1 + v$. Si la regresión auxiliar de la variable omitida es $X_2 = 1.5 + 0.5 X_1$, ¿a qué valor convergerá el estimador de la pendiente $\tilde{\beta}_1$?}
\begin{opts}
    \item $\tilde{\beta}_1$ converge a 4.0
    \item $\tilde{\beta}_1$ converge a 2.0
    \item $\tilde{\beta}_1$ converge a 5.0
    \item $\tilde{\beta}_1$ converge a 6.0
\end{opts}

\q{En una regresión simple con regresor endógeno $X$, se utiliza un instrumento exógeno $Z$. Se dispone de: $\widehat{Cov}(Z, Y) = 150$ y $\widehat{Cov}(Z, X) = 75$. ¿Cuál es el estimador de Variables Instrumentales $\hat{\beta}_{IV}$?}
\begin{opts}
    \item $\hat{\beta}_{IV} = 1.50$
    \item $\hat{\beta}_{IV} = 2.00$
    \item $\hat{\beta}_{IV} = 11250$
    \item $\hat{\beta}_{IV} = 0.50$
\end{opts}

\q{En la estimación por Mínimos Cuadrados en Dos Etapas (MC2E / 2SLS) con un regresor endógeno, la regresión de primera etapa arrojó un estadístico $F$ sobre los instrumentos excluidos de $F = 4.2$. Siguiendo la regla de oro de Stock, Wright y Yogo (2002) para detectar 'Instrumentos Débiles':}
\begin{opts}
    \item los instrumentos son débiles ($F < 10$), sesgando el estimador 2SLS hacia MCO.
    \item los instrumentos superan el umbral y garantizan inferencia no sesgada.
    \item la prueba de Durbin-Watson rechaza la existencia de endogeneidad.
    \item el modelo es exactamente identificado y debe estimarse por MCO.
\end{opts}

\q{Para verificar si una variable sospechosa $X_1$ es realmente endógena mediante la prueba de especificación de Durbin-Wu-Hausman en forma de regresión:}
\begin{opts}
    \item se añaden los residuos de la primera etapa $\hat{v}$ en la ecuación MCO y se evalúa su prueba $t$.
    \item se calcula la correlación simple entre $X_1$ y la constante tras aplicar primeras diferencias al panel.
    \item se reemplaza $X_1$ por su inversa y se contrasta si el intercepto de la nueva ecuación es nulo al $5\%$.
    \item se estima el modelo por MCO y por WLS y se exige que la diferencia entre ambos $R^2$ supere el $90\%$.
\end{opts}

\q{En un modelo con 1 variable endógena y 3 instrumentos excluidos (modelo sobreidentificado con $m = 3$ y $k_1 = 1$, grados de sobreidentificación $m - k_1 = 2$), la prueba de restricciones de sobreidentificación de Sargan arroja un estadístico de $Sargan = n R_{aux}^2 = 1.20$. Si el valor crítico al $5\%$ de $\chi^2_2$ es $5.99$, la conclusión es:}
\begin{opts}
    \item el modelo está subidentificado y requiere añadir más variables endógenas.
    \item la prueba indica que los instrumentos son débiles en la primera etapa.
    \item no se rechaza $H_0$: los instrumentos cumplen la restricción de exclusión.
    \item se rechaza $H_0$: los instrumentos son inválidos y correlacionados con $u$.
\end{opts}

\q{Al modelar diferencias de costos operativos entre regiones con 4 categorías (Costa, Sierra, Selva y Lima Metropolitana), si un economista incluye una variable ficticia (dummy) para cada una de las 4 regiones y además incluye el término constante en la regresión MCO:}
\begin{opts}
    \item se incurre en la Trampa de Variables Ficticias, generando multicolinealidad perfecta.
    \item el modelo gana máxima precisión al capturar efectos fijos completos de las cuatro zonas.
    \item los coeficientes miden directamente la desviación estándar de cada categoría regional.
    \item el $R^2$ ajustado aumenta artificialmente al cuadruplicar los grados de libertad.
\end{opts}

\q{En un modelo salarial $\ln(W) = \beta_0 + \beta_1 Educ + \beta_2 Mujer + \beta_3 (Educ \times Mujer) + u$, donde $Mujer$ es una dummy (1 si es mujer, 0 si es hombre), el retorno porcentual a un año adicional de educación para una mujer es:}
\begin{opts}
    \item $\beta_1$
    \item $\beta_1 \times \beta_3$
    \item $\beta_2 + \beta_3$
    \item $\beta_1 + \beta_3$
\end{opts}

\q{En una regresión lineal con variables estandarizadas (donde cada variable se transforma restando su media y dividiendo entre su desviación estándar $Z_X = \frac{X - \bar{X}}{S_X}$):}
\begin{opts}
    \item la varianza residual estimada es exactamente igual a $n - k$.
    \item el $R^2$ se convierte en un número estrictamente negativo.
    \item el intercepto estimado $\hat{\beta}_0$ es idénticamente cero.
    \item los errores estándar de todas las pendientes se duplican.
\end{opts}

\q{La matriz de proyección o 'Matriz Sombrero' (Hat Matrix) definida como $H = X(X'X)^{-1}X'$ cumple las propiedades algebraicas de ser:}
\begin{opts}
    \item asimétrica y no invertible bajo cualquier rango de $X$.
    \item simétrica ($H' = H$) e idempotente ($H^2 = H$).
    \item antisimétrica, definida negativa y ortogonal a $X$.
    \item estrictamente diagonal con ceros en la diagonal principal.
\end{opts}

\q{Una propiedad algebraica exacta de los residuos de MCO ($e = Y - X\hat{\beta}$) en cualquier modelo que incluya término constante es que:}
\begin{opts}
    \item la media de los residuos coincide con la varianza incondicional de $Y$.
    \item el producto de los residuos individuales es siempre igual a la unidad ($1$).
    \item la suma de los residuos es cero ($\sum e_i = 0$) y se cumple $X'e = 0$.
    \item los residuos son estocásticamente independientes entre sí para todo tamaño $n$.
\end{opts}

\q{La descomposición clásica de la varianza en MCO establece que la Suma Total de Cuadrados ($STC$) se descompone en Suma de Cuadrados Explicada ($SEC$) más Suma de Cuadrados Residuales ($SCR$): $STC = SEC + SCR$. Esta igualdad se cumple estrictamente siempre que:}
\begin{opts}
    \item el modelo lineal incluya un término constante (intercepto).
    \item todos los regresores sigan una distribución normal bivariada.
    \item el tamaño muestral exceda las 1000 observaciones continuas.
    \item no exista correlación lineal alguna entre $Y$ y $X$.
\end{opts}

\q{Para un coeficiente estimado $\hat{\beta}_1 = 2.50$ con $EE(\hat{\beta}_1) = 0.25$ en una muestra con $n - k = 60$ grados de libertad (valor crítico $t_{0.025} = 2.00$), el intervalo de confianza al $95\%$ es:}
\begin{opts}
    \item $[2.45, \; 2.55]$
    \item $[2.25, \; 2.75]$
    \item $[2.00, \; 3.00]$
    \item $[1.50, \; 3.50]$
\end{opts}

\q{Al realizar una predicción individual puntual de $Y_0$ para un vector de valores de regresores $x_0$, la varianza del error de predicción individual $Var(Y_0 - \hat{Y}_0)$ es:}
\begin{opts}
    \item $\sigma^2 [x_0' (X'X)^{-1} x_0]$
    \item $\sigma^2 [n + x_0' (X'X) x_0]$
    \item $\sigma^2 [1 + x_0' (X'X)^{-1} x_0]$
    \item $\sigma^2 [1 - x_0' (X'X)^{-1} x_0]$
\end{opts}

\q{En la selección entre especificaciones de modelos rivales, el Criterio de Información de Schwarz (BIC / SIC) se diferencia del Criterio de Akaike (AIC) en que:}
\begin{opts}
    \item el AIC solo se utiliza en modelos de variables instrumentales y panel dinámico.
    \item el BIC penaliza más parámetros adicionales ($\ln(n)k$ vs $2k$ en AIC).
    \item el BIC exige que el $R^2$ ajustado supere $0.80$ para seleccionar modelos.
    \item el BIC no incorpora ningún término de penalización por grados de libertad.
\end{opts}

\q{La prueba de normalidad de Jarque-Bera sobre los residuos de MCO evalúa si el coeficiente de asimetría muestral $S$ y la curtosis muestral $K$ coinciden con los de una distribución normal teórica ($S=0, K=3$). Su estadístico es:}
\begin{opts}
    \item $JB = \frac{n}{24} \left( S^2 + \frac{K^2}{4} \right) \sim \chi^2_4$
    \item $JB = \frac{n}{6} \left( S^2 + (K - 3)^2 \right) \sim \chi^2_2$
    \item $JB = \frac{n}{12} \left( (S - 1)^2 + \frac{(K - 3)^2}{2} \right) \sim F_{2, n}$
    \item $JB = \frac{n}{6} \left( S^2 + \frac{(K - 3)^2}{4} \right) \sim \chi^2_2$
\end{opts}

\q{La prueba RESET de Ramsey (Regression Equation Specification Error Test) se utiliza en econometría para:}
\begin{opts}
    \item determinar el número óptimo de rezagos autorregresivos $p$ y de media móvil $q$ en series ARIMA.
    \item detectar errores de forma funcional añadiendo potencias $\hat{Y}^2, \hat{Y}^3$ con test $F$.
    \item contrastar si existe heterocedasticidad condicional autorregresiva en los residuos de la muestra.
    \item calcular la elasticidad precio cruzada en sistemas de ecuaciones aparentemente no relacionadas.
\end{opts}

\q{Al estimar un modelo de MCO con restricciones lineales $R\beta = r$ (por ejemplo, rendimientos constantes a escala en una función Cobb-Douglas $\beta_1 + \beta_2 = 1$), el estimador MCO restringido $\hat{\beta}_R$ frente al no restringido $\hat{\beta}_{NR}$:}
\begin{opts}
    \item impide calcular matriz de covarianzas o errores estándar robustos de White.
    \item siempre presenta $SCR_R \le SCR_{NR}$, logrando un mayor ajuste de $R^2$.
    \item se vuelve inconsistente asintóticamente aunque las restricciones sean ciertas.
    \item siempre presenta $SCR_R \ge SCR_{NR}$, al optimizar en un espacio acotado.
\end{opts}

